%% The following is a directive for TeXShop to indicate the main file
%%!TEX root = diss.tex

%=======================================
\chapter{Two-dimensional Random Walk Models}
\label{ch:2DWalks}
%=======================================
%=======================================
\section{Introduction}
%=======================================
The power of the operator notation really starts to show here as it allows us to easily derive functional equations for the two-dimensional problem.
\begin{definition} Let $Q = Q(x_1,x_2,a_1,a_2;t)$ be the generating function enumerating two-dimensional quarter-plane random walks with interactions with the axes, where $t$ counts the length, the $x$ monomials encode the endpoints and for each $k$, $a_k$ counts the interactions with the $x_k = 0$ line.
\end{definition}
The derivation of the functional equation is done in the same way as in the one-dimensional case. Starting with $Q$, extend by one step via $tS$, apply operators $L$ and $E$ and obtain $Q-1$. That is,
\begin{align}
	(1-tLES)Q &= 1\,,\\
	(L^{\um 1} - tES)Q &= L^{\um 1}(1) = \alpha\,,
\end{align}
where $L,\,E$ and $\alpha$ are to be determined in order to find the master functional equation. For $E$ it is enough to consider where walks may end when we add a step. The furthest they can reach are coordinates $(m,n)$ where one of them is equal to $\um 1$ and the other is non-negative. It is enough to first remove walks with monomial $\ol{x}_2x_1^k$ and then walks with monomial $\ol{x}_1x_2^k$ for arbitrary $k$. Each of these actions can be performed via an operators $E_2$ and $E_1$ respectively, where
\begin{align}
	E_k = 1 - \ext{\ol{x}_k}\,.
\end{align}
This allows us to write
\begin{align}
	E = E_1 E_2 = 1-\ext{\ol{x}_1}-\ext{\ol{x}_2}+\ext{\ol{x}_1\ol{x}_2} = E_2 E_1\,,
\end{align}
which also shows that $E_1$ and $E_2$ commute. It's interesting to note that this multiplication of operators captures the inclusion-exclusion nature of the problem. We remove all walks with monomials containing $\ol{x}_1$ or $\ol{x_2}$ but we've removed some of the twice, namely those with monomial $\ol{x}_1\ol{x}_2$, and so we need to add those back in.

Similarly to operators $E_k$ we define $L_k = 1+ b_k[x_k^0]$ with inverses $L_k^{\um 1} = 1 - c_k[x_k^0]$. We find that $L = L_1L_2 = L_2L_1$ just like $E$. The reason that $L$ works as advertised is that each $L_k$ acts by multiplying monomials by $a_k$ and that the multiplicative action is commutative. It doesn't matter if we apply $L_1$ first and then $L_2$ or some fancier operator that does the marking in one go. To find $\alpha$ we just note that $L_k^{\um 1}a = (1-c_k)a$ for a constant $a$ and thus
\begin{align}
	L^{\um 1}(1) = (1-c_1)(1-c_2) = \ol{a}_1\ol{a}_2\,.
\end{align}
Next we'll find the master functional equation. We have
\begin{align}
	L^{\um 1} Q = (1-c_1[x_1^0])(1-c_2[x_2^0]) Q = Q - c_1 Q_1^0 - c_2 Q_2^0 + c_1c_2 \Qf{0,0}
\end{align}
and
\begin{align}
\begin{split}
	ESQ  &= (1-\ext{\ol{x}_1})(1-\ext{\ol{x}_2})SQ \\
	&= SQ - \ol{x}_1S_1^{\um 1}Q_1^0 - \ol{x}_2S_2^{\um 1}Q_2^0 + \ol{x}_1\ol{x}_2\Sf{\um 1,\um 1}\Qf{0,0}
\end{split}
\end{align}
where in the calculation we've dropped terms like $Q_i^j$ for $j <0$ since the model is restricted to the quarter plane. Combining all the terms computed so far leads to the master functional equation
\begin{align*}
	(1-tS)Q - (c_1-t\ol{x}_1S_1^{\um 1}) Q_1^0- (c_2-t\ol{x}_2S_2^{\um 1}) Q_2^0+ (c_1c_2-t\ol{x}_1\ol{x}_2S_{1,2}^{\um 1,\um 1})Q_{1,2}^{0,0}  = \ol{a}_1\ol{a}_2\,.
\end{align*}
We note the pattern in the coefficients with the following definition.
\begin{definition} For $I \subseteq [d]$ let
\begin{align}
K_I = c_I - t\ol{x}_I S_I^{\um 1}
\end{align}
where if $I = \emptyset$, $K_I = K$.
\end{definition}
This allows us to re-write the equation neatly.
\begin{corollary} The generating function function $Q$ satisfies the master functional equation
\begin{align}
	KQ - K_1 Q_1^0 - K_2 Q_2^0 + K_{1,2} \Qf{0,0} = \ol{a}_1\ol{a}_2\,.
\end{align}
\end{corollary}
We can also perform coefficient extractions like in the one-dimensional case and obtain further functional equations. The derivation is quite mechanical and very similar to the one-dimensional case so we will omit it here.
\begin{corollary} The generating function $Q$ satisfies the additional functional equations
\begin{align*}
(1-ta_1 S_1^0)Q_1^0 - (c_2-ta_1\ol{x}_2S_{1,2}^{0,\um 1})Q_{1,2}^{0,0} - a_1 t(S_1^{\um 1}Q_1^1-\ol{x}_2S_{1,2}^{\um 1,\um 1}Q_{1,2}^{1,0}) &=\ol{a}_2\,,\\
(1-ta_2 S_2^0)Q_2^0 - (c_1-ta_2\ol{x}_1S_{1,2}^{\um 1,0})Q_{1,2}^{0,0} -a_2t(S_2^{\um 1}Q_2^1-\ol{x}_1S_{1,2}^{\um 1,\um 1}Q_{1,2}^{0,1}) &=\ol{a}_1\,,\\
(1-a_1 a_2 t \Sf{0,0})Q_{1,2}^{0,0} - ta_1a_2t(\Sf{\um 1,\um 1}\Qf{1,1}+\Sf{\um 1,0}\Qf{1,0}+\Sf{0,\um 1}\Qf{0,1})&= 1
\end{align*}
where $\Sf{0,0} = 1$ if we have resting steps and zero otherwise.
\end{corollary}

We should contrast this quick analysis to the one done without the operator notation. We would have to use inclusion-exclusion to determine which terms are being kept and which terms are being discarded, and then we would have to take care of double-counting. The procedure can be somewhat complicated and prone to error. Furthermore, the complexity of the equation increases exponentially with the number of extractions necessary. The operator notation allows us to derive the functional equatiosn in an elegant way that is rather easy to generalize. It also allows us to encapsulate away the unnecessary details in the marked cases, so that we can think of inverting operators, rather than do coefficient extraction and back substitution (which would be the original way). Later we'll cover the general case in more detail.

%=======================================
\section{Trivial First Examples}
%=======================================

Quarter-plane models with small steps and without interactions have been investigated extensively but the case with interaction has not been investigated as much and we are interested in finding the generating function $Q$ explicitly, or at least $\Qf{0,0}$. This appears to be difficult in most nontrivial cases so in general we may have to resort to finding if $Q$ or $\Qf{0,0}$ is D-finite in the $t$ variable. Apart from that recent research into two-dimensional random walks with interactions has shown that there are models where we have an explicit relationship between $Q$ or $\Qf{0,0}$ and unmarked variants. As we saw in the previous chapter the on-dimensional cases admit a linear relationship with polynomial coefficients. There two-dimensional models that exhibit a rational relationship. It is possible that other models don't even have such a relationship with their unmarked counterparts.

Our work can be split into several cases of increasing difficulty. We start with trivial models, continue through some semi-trivial models, go through some nontrivial models and finish at some difficult problems in the latter chapter. The trivial models are those where we don't have to do any work to find the generating functions. There are models with walks that never leave the origin. However complicated $\mc S$ it doesn't matter because $Q = 1$ when we don't have resting steps and $1/(1-a_1a_2S_0t)$ when we have a resting step. There are also the semi-triviall models where walks leave the origin but never come back to any axis. We'll solve such models explicitly later on but for now it suffices to say that $\Qf{0,0} = 1$ in the case with no resting steps and $\Qf{0,0} = 1/(1-a_1a_2S_0t)$ in the case with resting steps. In each of these situations we have a relationship between the interacting and noninteracting models. In any of these simple cases we can write $B\Qf{0,0} = 0$ and sometimes $BQ = 0$ where $B = B_1 B_2$ and $B_k = 1-[b_k^0]$ as in the one-dimensional case.

Here are all the easy cases of step sets that lead to walks that either don't leave the origin or leave it and never come back. We're not allowing resting steps here. There are 80 such sets among the 256 non-resting possible cases. Of course we could use symmetries to eliminate cases bu we won't do that.

There are $2^5 = 32$ types of walks that don't leave the origin. Those are walks with steps among the set $\{$NW,W,SW,S,SE$\}$. Next we look at the walks that leave the origin but never come back.  There is the case of $S = \{NE\}$ that never comes back to any axis. Let's look at those walks that never come back to the $x_1$ axis. We have some steps that go in the positive $x_2$ direction, that is one of $\{$N, NE, E$\}$ together with some steps from $\{$NW, W, SW$\}$. If we were to include any other steps we would come back to the $x_1$ axis. If we only have a N step, we have 8 new cases, if we have only a NE step we have 3 new cases and if we have both, we have 4 new cases. This gives a total of 15 new cases. There's an equivalent number of walks that leave the $x_2$ axis and never come back to it. The running total is $32 + 1  + 15 + 15 = 63$.

To finish off we have 17 cases of walks that leave the origin, eventually interact with both axes but never come back to the origin. For this we must only include steps in such a way that we cannot pick a selection of steps that will decrease both the $x_1$ and $x_2$ coordinates. We have to pick a subset of steps from $\{$N, NE, E$\}$ so the only other steps allowed are NW and SE. The new cases coming from $\{$N, NE, E$\}$, $\{$N,E$\}$, $\{$N, NE $\}$, $\{$NE, E$\}$, $\{$N$\}$, $\{$NE$\}$, $\{$E$\}$ are 4, 4, 2, 2, 1, 2, 2 respectively.

Before we consider walks that come back to the origin we'll look at the behavior $Q$ depending on whether $S$ factors in some nice way or not.

%=======================================
\section{Simple case: Factoring $S$ into a product}
%=======================================

We take a moment to mention a result that falls out of the operator notation and which was covered in more generality by Steve et al. 
\begin{example}
We have a special case when $S$ factors as $S = S_1(x_1)S_2(x_2)$. We have
\begin{align}
	 (\Gamma S)^k = (\Gamma_1\Gamma_2S_1S_2)^k = (\Gamma_1S_1)^k(\Gamma_2S_2)^k
\end{align}
since the $\Gamma _1$ and $\Gamma _2$ operators commute and the $\Gamma _iS_i$ operators also commute since they deal with different variables. The functional equation for $Q$ is
\begin{align}
	(1-t\Gamma S)Q = 1\,.
\end{align}
We can compute $Q$ directly via 
\begin{align}
	Q = \sum_{k = 0}^{\infty} t^k (\Gamma S)^k = \sum_{k = 0}^{\infty} t^k (\Gamma_1S_1)^k(\Gamma_2S_2)^k = Q_1(x_1, a_1,t) \odot_t Q_2(x_2,a_2,t)
\end{align}
where $\odot_t$ denotes the Hadamard product of two functions with respect to variable $t$. From now on we'll drop the subscript $t$ unless it is needed. Here $Q_i$ is given by
\begin{align}
	Q_ = \sum_{k = 0}^{\infty} t^k (\Gamma_iS_i)^k
\end{align}
and by reversing the iteration we obtain the functional equation
\begin{align}
	(1-t\Gamma_iS_i)Q_i = 1\,.
\end{align}
In other words, $Q_i(x_i,a_i,t)$ is the generating function for a one-dimensional problem with Laurent step polynomial $S_i(x_i)$. In the case of a factoring of this sort, we are able to glue the two-dimensional solution out of two one-dimensional solutions. (Of course this generalizes easily to many dimensions and gives an infinite family of solutions to higher-dimensional problems where $S$ factors into lower dimensional $S_i$'s). If we want to look at the generating function for quarter plane walks of length $k$ it is enough to look at the product of two generating functions for one-dimensional walks of length $k$.
\end{example}
\begin{example}
There is another way to see the Hadamard product solution without the iteration method. Suppose that $Q_i$ is a solution to the functional equation $(L_i^{\um 1} - tE_iS_i)Q_i = \ol{a}_i$ for $i = 1$ and $i = 2$. Then we claim that $Q = Q_1 \odot Q_2$ is a solution to the functional equation $(L^{\um 1} - tES)Q = \ol{a}_1\ol{a}_2$. To this end we calculate
\begin{align}\begin{split}
	(L^{\um 1}-tES)Q &= L^{\um 1} (Q_1\odot Q_2) - tES(Q_1\odot Q_2)\\
	& = (L^{\um 1}_1Q_1\odot L^{\um 1}_2Q_2) - t((E_1S_1Q_1)\odot (E_2S_2Q_2))\,.
\end{split}\end{align}
Then we substitute $ L^{\um 1}_i Q_i = \ol{a}_i + tE_iS_iQ_i$ into the first term to obtain
\begin{align}\begin{split}
	 (L^{\um 1}_1Q_1\odot L^{\um 1}_2Q_2) &= (\ol{a}_1+ t E_1 S_1 Q_1)\odot(\ol{a}_2 + t E_2 S_2 Q_2)\\
	 & = \ol{a}_1\odot\ol{a}_2 + \ol{a}_1\odot (t E_2 S_2 Q_2) + \ol{a}_2\odot(t E_1 S_1 Q_1) \\
	 &\quad+ (t E_1 S_1 Q_1)\odot(t E_2 S_2 Q_2))\,.
\end{split}\end{align}
We have
\begin{align}
	\ol{a}_i\odot (t E_j S_j Q_j) = \ol{a}_i [t^0]t E_j S_j Q_j = 0
\end{align}
since we are dealing with power series in $t$. Therefore
\begin{align}
	(L^{\um 1}_1Q_1\odot L^{\um 1}_2Q_2) =  \ol{a}_1\ol{a}_2 +t ((E_1 S_1 Q_1)\odot (E_2 S_2 Q_2))
\end{align}
and upon substitution we are left with $\ol{a}_1\ol{a}_2$ which is what we needed in the first place.
\end{example}

\begin{lemma} Let $Q$ be the generating function for a two-dimensional random walk model with small steps and interactions on the quarter plane with $S = S_1(x_1)S_2(x_2)$. Let $Q_i$ be the generating function for a one-dimensional random walk model with small steps and interaction with the origin and step polynomial $S_i$. Then
\begin{align}
	Q(x_1,x_2,a_1,a_2;t) = Q_1(x_1,a_1;t)\odot Q_2(x_2,a_2;t)\,.
\end{align}
\end{lemma}

This factoring also generalizes to higher dimensions. Even though we haven't shown the higher-dimensional functional equation yet its form is not difficult to guess.
\begin{theorem} Consider a $d$-dimensional random walk model with small steps restricted to the region with $x_k \ge 0$ and having interactions with the $x_k = 0$ hyperplanes. Suppose that $S(\row xd) = S_1\ldots S_r$ where $S_i$ and $S_j$ don't share variables when $i\neq j$. Then the generating function $Q$ is given by
\begin{align}
	Q = Q_1 \odot \cdots \odot Q_r
\end{align}
where $Q_k$ is the sub $d$-dimensional problem with smalls steps restricted to the region with $x_i \ge 0$ and having interactions with the $x_i = 0$ hyperplane where $x_i$ is a varible present in $S_k$.
\end{theorem}

We may glue lower-dimensional solutions to obtain solutions to higher-dimensional problems. In particular we've just shown the result that we can solve an infinite family of nontrivial random walk models.

\begin{example} Consider the two-dimensional model with $S = (x_1+\ol{x}_1)(x_2+\ol{x}_2)$. The one dimensional problems are the same and have solutions
\begin{align}
	\widetilde{Q}_0 = \frac{\sigma}{at-b\sigma}\com \widetilde{Q} = \frac{t(x-\sigma)}{x(1-t(x+\ol{x})(at-b\sigma)}
\end{align}
where
\begin{align}
	\sigma = \frac{1-\sqrt{1-4t^2}}{2t}\,.
\end{align}
We can expand and simplify $\widetilde{Q}_0$ and $\widetilde{Q}$ to find
\begin{align}
	\widetilde{Q}_0 &= \frac{2-a-a\sqrt{1-4t^2}}{2(1-a+a^2t^2)}\,\\
	\widetilde{Q} &= \frac{((a-2)t+(1-a)x+2axt^2)+(at+(1-a)x)\sqrt{1-4t^2}}{2(1-a+a^2t^2)(x-t-tx^2)}\,.
\end{align}
We have the expansion
\begin{align}
	\widetilde{Q} = 1 + xt + (a+x^2)t^2+x(1+a+x^2)t^3 + \cdots\,.
\end{align}
This allows us to write
\begin{align}
\begin{split}
	Q(x_1,x_2,a_1,a_2;t) =& \widetilde{Q}(x_1,a_1;t) \odot \widetilde{Q}(x_1,a_2;t)\\
		=&1 + x_1x_2t + (a_1+x_1^2)(a_2+x_2^2)t^2\\
		&+x_1x_2(1+a_1+x_1^2)(1+a_2+x_2^2)t^3 + \cdots\,.
\end{split}
\end{align}
However obtaining an explicit function via the Hadamard product seems to be rather messy. The function will not be algebraic since the unmarked variant is not algebraic. But at least the function $Q$ is D-finite since both functions in the product are.
\end{example}

%=======================================
\section{Kernel Group Considerations}
%=======================================
It is important to find more functional equations if we can. One way to do this is by applying a transformation on the pair $(x_1,x_2)$ that leaves as many parts of the functional equation as we can but still adding some information in the form of a new equation. One way to do this is to cosider mapping $x_1$ and $x_2$ to some quantities separately and then consider the full group generated by such maps. Since we usually eliminate the kernel in front of $Q$ it is useful to consider transformations that preserve the kernel.  We'll call this group $G(K)$, although it won't be the full group of rational transformations preserving the kernel. For example $(x_1,x_2) \mapsto (x_2,x_1)$ will not be in the group. Any transformation on $(x_1,x_2)$ that preserves the kernel $K = 1 -tS$ must necessarily preserve the step polynomial $S$. Here
\begin{align}
	S = \alpha \ol{x} + \beta + \gamma x
\end{align}
where $x = x_1$ or $x_2$ and $\alpha,\beta,\gamma$ change correspondingly. We consider sending $x \mapsto y$ and solve for $y$ to determine the possible transformations we can perform. That is we have
\begin{align}
	\alpha \ol{x} + \gamma x &=\alpha \ol{y} + \gamma X\\
	(\alpha - \gamma xy)(y-x) &= 0\,.
\end{align}
We have cases
\begin{itemize}
\item $\alpha = 0, \gamma = 0$. The equation is then trivially true and any transformation preserves the kernel although this is not very useful though.
\item $\alpha = 0, \gamma \neq 0$. Then $ y = x$ and only the identity transformation preserves the kernel.
\item $\alpha \neq 0, \gamma = 0$. Here $y = x$ as well.
\item $\alpha \neq 0, \gamma \neq 0$. Here we either have $y = x$ or $X\neq x$ and $y = \ol{x}\alpha/\gamma = \ol{x} S_{\um 1}/S_1$.
\end{itemize}
In conclusion, either every transformation works and we don't have an interesting case or we have the identity transformation plus a possible second one. The last transformation creates a nontrivial group and we can have groups of various orders. Let
\begin{align}
	g_1 : x_1 \mapsto \ol{x}_1\frac{S_1^{\um 1}}{S_1^{1}},\quad g_2 : x_2 \mapsto \ol{x}_2\frac{S_2^{\um 1}}{S_2^{1}}
\end{align}
whenever the $S$ quantities are non-zero. Note that $g_i(g_i(x_1)) = x_1$, that is they are involutions. We have three cases for each $x$
\begin{itemize}
\item 1: Any transformation can be applied.
\item 2: We only have the identity transformation.
\item 3: We have two transformations.
 \end{itemize}
 We'll combine cases for analysis in the form "case $i.j$" to mean we have case $i$ for $x_1$ and case $j$ for $x_2$.

\begin{itemize}
\item [Case 1.1:] Here we can apply any transformation on both $x_1$ and $x_2$. From the transformation analysis this means that $S_1^{\um 1} = S_1^1 = S_2^{\um 1} = S_2^1 = 0$ and thus means that we have no steps that go away from the origin. Therefore $S = \Sf{0,0}$. Thus we have only two models - no steps or a resting step. 
\item[Case 1.2:] Here any transformation applies on $x_1$, so $S_1^{\um 1} = S_1^1 = 0$ so the model has no steps that step away from the origin in the $x_1$-direction. Since the only transformation on $x_2$ is the identity transformation we either have $S_2^{\um 1} = 0$ or $S_2^1 = 0$ in which case the model consists of only an up step or a down step, together with a possible resting step. Similarly, the 2.1 case involves a single $x_1$ or $\ol{x}_1$ step with a possible resting step.
\item[Case 1.3:] (and 3.1): This case is similar to the previous one except that now we have two transformations for $x_2$ and this means $S_2^{\um 1}, S_2^1 \neq 0$. This means we must have either both vertical steps $x_2$ and $\ol{x}_2$ in the 1.3 case or two horizontal steps $x_1$ and $\ol{x}_1$ in the 3.1 case.
\item[Case 2.2:] In this case both transformations are the identity transformations. Therefore exactly one of $S_1^{\um 1}$ and $S_1^1$ is zero and exactly one of $S_2^{\um 1}$ and $S_2^1$ is zero. WLOG, assume that $S_1^1 \neq 0, S_2^1 \neq 0$ (the other models can be obtained by a rotation around the origing of 90, 180 and 270 degrees). The possible models are (I should include the drawings): $\{$N,E,NE$\}$, $\{$N,NE$\}$, $\{$E,NE$\}$, $\{$N,E$\}$, $\{$NE$\}$. In other words, the subsets of $\{$N,E,NE$\}$ except the ones going either only in the N direction or only in the E direction. This is a case of group $G(K)$ or order 1.
\item[Case 2.3:] In this case there is only one transformation for $x_1$ and two for $x_2$. This means $S_1^{\um 1} = 0$ (positive case) or $S_1^{1} = 0$ (negative case). Also $S_1^{\um 1}, S_2^1 \neq 0$ so there must be at least one step going in the positive $x_2$ direction and one in the negative $x_2$ direction. We have 17 cases (and another 17 including a resting step). Of course in the 3.2 case we have the meaning of $x_1$ and $x_2$ reversed.This is a case of group $G(K)$ of order 2.
\item[Case 3.3:] This is the most interesting case with two transformations for $x_1$ and $x_2$ each. This means that there must be at least one step going in a positive and one going in a negative coordinate for each of $x_1$ and $x_2$. There are a lot of such cases (in the hundreds). Previous results show that the group $G(K)$ has orders 4, 6, 8 and infinity depending on the step set in question.
 \end{itemize}
 
 Later we'll see that these cases correspond to the idea of dimensionality of the model. Case 1.1 is zero-dimensional. Cases 1.2, 1.3 (and 3.1) are one-dimensional. Cases 2.2, 2.3 and 3.3 are two-dimensional. It should not be a suprise that the 0 and 1-dimensional cases will be easy while the two-dimensional cases will become progressively more difficult.
 
 %=======================================
\section{Solving basic 2D models}
%=======================================

%=======================================
\subsection{Cases 1.1, 1.2, 1.3}
%=======================================
A quick glance at the step models involved shows that we are dealing with essentially one-dimensional problems. This means they can be solved easily explicitly. One can also note that these models are all factorable
\begin{align}
	S(x_1,x_2) = 1\cdot S(x_2)\,.
\end{align}
Both $Q_1$ and $Q_2$ are functions we solved for in the one-dimensional case
\begin{align}
	Q_1(x_1,a_1,t) = \frac{1}{1-a_1 t},\quad Q_2(x_2,a_2,t) = q(x_2,a_2,t)
\end{align}
for some solution $q$ from that section. Therefore
\begin{align}\begin{split}
	Q(x_1,x_2,a_1,a_2,t) &= Q_1(x_1,a_1,t) \odot Q_2(x_2,a_2,t) \\
	&= \frac{1}{1-a_1 t}\odot q(x_2,a_2,t) = q(x_2,a_2,a_1 t)\,.
\end{split}\end{align}
Similarly in the 2.1, 3.1 cases
\begin{align}\begin{split}
	Q(x_1,x_2,a_1,a_2,t) &= Q_1(x_1,a_1,t) \odot Q_2(x_2,a_2,t) \\
	& =  q(x_1,a_1,t) \odot \frac{1}{1-a_2 t}= q(x_1,a_1,a_2 t)
\end{split}\end{align}
and this completes the solution. What about the relationship with the $B_1B_2$ transformation? If $f(a,t)$ is the function that goes together with $q(x,a,t)$ then the function $F(a_1,a_2,t)$ that goes together with $Q(x_1,x_2,a_1,a_2,t)$ is given by
\begin{align}
	F(a_1,a_2,t) = f(a_2,a_1 t)\,.
\end{align}
Then we have $B_2 (FQ) = 0$ and $B_1B_2(FQ) = 0$.

%=======================================
\subsection{Case 2.2 - group of order 1}
%=======================================
A quick glance at the step models reveals that this case is mostly trivial. We have a lot of models that never leave the origin and those have a trivial solution
\begin{align}
	Q(x_1,x_2,a_1,a_2) = \frac{1}{1-a_1a_2 \Sf{0,0} t}\,.
\end{align}
We also have a lot of models that are essentially one-dimensional with a single step going along an axis and some steps going off-axis where the latter steps can never be taken by a random walk in this model. They are essentially cases already covered in the previous section. Suppose that the single step is taken in the positive $x_1$ direction. This is equivalent to the $s_{-1} = 0$ case in 1D with function
\begin{align}
	q(x,a,t) = \frac{1-s_0 t}{(1-st)(1-a s_0 t)} = \frac{1-s_0 t}{(1-x_1s_1 t)(1-a s_0 t)}\,.
\end{align}
Then we have
\begin{align}\begin{split}
	Q(x_1,x_2,a_1,a_2,t) &= Q_1(x_1,a_1,t) \odot Q_2(x_2,a_2,t) \\
	&=  q(x_1,a_1,t) \odot \frac{1}{1-a_2 t}= q(x_1,a_1,a_2 t)
\end{split}\end{align}
and therefore
\begin{align}
	Q(x_1,x_2,a_1,a_2,t) = \frac{1-a_2 \Sf{0,0} t}{(1-(x_1+\Sf{0,0})a_2 t)(1-a_1a_2 \Sf{0,0}t)}\,.
\end{align}
The interesting models in case 2.2 are those with non-negative steps that we listed. Here
\begin{align}
	S = \alpha x_1 + \beta x_2 + \gamma x_1 x_2 + \delta\,.
\end{align}
The functional equation and its zero extractions are
\begin{align}\begin{split}
(1-tS)Q - c_1 Q_1^0- c_2 Q_2^0+ c_1c_2 Q_{1,2}^{0,0}  &= \ol{a}_1\ol{a}_2\,,\\
(1-ta_1 (\delta+\beta x_2))Q_1^0 - c_2Q_{1,2}^{0,0} &=\ol{a}_2\,,\\
(1-ta_2 (\delta+\alpha x_1))Q_2^0 - c_1Q_{1,2}^{0,0}  &=\ol{a}_1\,,\\
(1-a_1 a_2 t \delta)Q_{1,2}^{0,0} &= 1\,.
\end{split}\end{align}
We solve the problem by a simple back-substitution of the variables.
\begin{align}
\Qf{0,0} & = \frac{1}{1-a_1a_2\delta t}\,,\\
Q_1^0 & = \frac{\ol{a}_2}{1-t a_2 S_1^0}\left(1+\frac{b_2}{1-a_1a_2 \delta t}\right) = \frac{1 -  a_1 \delta t}{(1-t a_1S_1^0)(1-a_1a_2\delta t)}\,,\\
Q_2^0 & = \frac{\ol{a}_1}{1-t a_1 S_2^0}\left(1+\frac{b_1}{1-a_1a_2 \delta t}\right) = \frac{1 -  a_2 \delta t}{(1-t a_2S_2^0)(1-a_1a_2\delta t)}\\
Q &= \frac{1}{1-tS}(\ol{a}_1\ol{a}_2 + c_1 Q_1^0 + c_2 Q_2^0 - c_1 c_2 Q_{1,2}^{0,0})\\
&=\frac{1}{1-tS}\left (\ol{a}_1\ol{a}_2 + \frac{1}{1-ta_1a_2 \delta}\left(c_1 \frac{1 -  a_1 \delta t}{1-t a_1S_1^0} + c_2 \frac{1 -  a_2 \delta t}{1-t a_2S_1^0} - c_1 c_2 \right) \right)\\
&= \frac{1}{1-tS}\left(\ol{a}_1\ol{a}_2 + \frac{1}{1-a_1a_2\delta t}(c_1 f_1(a_1,t) + c_2 f_2(a_2,t)-c_1c_2)\right)
\end{align}
where the last line explicitly shows that $Q$ is composed of functions that depend only on one of $a_1$, $a_2$. If we want to find a relation of the type $B(FQ) = 0$ then we must first find a suitable function $F$. First of all the $1-tS$ part doesn't play a role since it doesn't depend on $a_1$ or $a_2$. Then later the variables $c_1$ and $c_2$ will vanish when extracting the coefficients of $b_1^0$ and $b_2^0$. We'll now show that the function
\begin{align}
	F(a_1,a_2,t) = 1 - a_1 a_2 \delta t
\end{align}
works exactly. First
\begin{align}
B((1 - a_1 a_2 \delta t)\ol{a}_1\ol{a}_2) = B(\ol{a}_1\ol{a}_2 - t \delta) = B(\ol{a}_1\ol{a}_2) - B(t \delta) = \ol{a}_1\ol{a}_2- \ol{a}_2-\ol{a}_1 + 1 = c_1 c_2\,.
\end{align}
Then
\begin{align}
	B(c_1 f_1+ c_2 f_2-c_1c_2) = c_1 f_1 + c_2 f_2 -c_1c_2 - c_2 f_2 - c_1 f_1 = -c_2 c_2\,.
\end{align}
Putting the two quantities together along with dividing by $1-tS$ shows that $B(FQ) = 0$. Of course this works for $G = \Qf{0,0}$ too.

%=======================================
\subsection{Case 2.3 - group of order 2}
%=======================================
There are two possibilities here - positive and negative. In the negative case $S_1^1 = 0$ and the steps have a non-positive $x_1$ component. Due to the fact that a lot of steps are not going to be taken by random walks in those models, all of these models reduce to the cases of $S$ consising of a rest step and maybe a single step in a positive direction. This takes care of 34 models (or 68 if we include a rest step). The interesting models are those with a positive $x_1$ component, or in other words, models where $S_1^{\um 1} = 0$. The calculations for these models become increasingly more complicated. The functional equation and its extracted forms read as
\begin{align}\begin{split}
(1-tS)Q - c_1 Q_1^0- (c_2 -t\ol{x}_2 S_1^{\um 1})Q_2^0+ c_1c_2 Q_{1,2}^{0,0}  &= \ol{a}_1\ol{a}_2\,,\\
(1-ta_1 S_1^0)Q_1^0 - (c_2 - ta_1 \ol{x}_2 \Sf{0,-1})Q_{1,2}^{0,0} &=\ol{a}_2\,,\\
(1-ta_2 S_2^0)Q_2^0 - c_1Q_{1,2}^{0,0}  - a_2 t S_1^{\um 1} Q_2^1&=\ol{a}_1\,,\\
(1-a_1 a_2 t \Sf{0,0})Q_{1,2}^{0,0} - ta_1 a_2 \Sf{0,-1} \Qf{0,1} &= 1\,.
\end{split}\end{align}
We can identify a subproblem that is essentially one-dimensional and therefore already solved. For this we look at the second and fourth equations with the substitution $a_i = a, c_i = c, x_2 = x, ta_1 = \tau, S_1^0 = T, Q_1^0 (x_2)= q(x)$. These equations become
\begin{align}
(1-\tau T)q - (c - \tau  \ol{x} T_{-1})q_0 &=\ol{a}\,,\\
(1-a \tau T_0)q_0 - a\tau \Sf{0,-1} q_1 &= 1
\end{align}
where $q_0 = q(a,\tau)$. This is the one-dimensional problem again. The solution here is $q(x,a,\tau)$ obtained by the methods in the previous chapter so the original problem has a solution
\begin{align}
	Q_1^0(x_2) = q(x_2,a_2,a_1t),\quad Q_{1,2}^{0,0} = q_0(a_2,a_1t)\,.
\end{align}
The reason the method worked in this case is that we have steps that keep a random walk staying on the $x_2$ axis and steps that take it away from the $x_2$ but no steps that have it return to it. This is similar to the case where a random walk leaves the $x_2$ axis never to come back. Essentially the problem splits into a one-dimensional one regarding the $x_2$ axis and a further 2D one once we step off the axis. We re-write the master equation as
\begin{align}
	(1-tS)Q - (c_2 -t\ol{x}_2 S_1^{\um 1})Q_2^0  &= \ol{a}_1(\ol{a}_2 + b_1 q  -b_1c_2q_0)  =\ol{a}_1\phi(x_2)\,.
\end{align}
Since $S_1^{\um 1},S_2^1 \neq 0$ we have a root $x_2 = \sigma_2$ that allows us to remove the $(1-tS)Q$ term and get
\begin{align}
	Q_2^0  &=-\frac{\ol{a}_1\phi(\sigma_2)}{ c_2 -t\ol{\sigma}_2 S_1^{\um 1}}
\end{align}
so this allows us to write $Q$ as
\begin{align}
	Q = \frac{1}{1-tS}\left(\ol{a}_1\phi(x_2) - \frac{c_2 -t\ol{x}_2 S_1^{\um 1}}{c_2 -t\ol{\sigma}_2 S_1^{\um 1}}\ol{a}_1\phi(\sigma_2)\right)\,.
\end{align}
Of course, this can be expanded and simplified based on what $q$ and $q_0$ are from the one-dimensional setting but the important thing is that we can obtain the solution in this case via one-dimensional equations.

%=======================================
\subsection{Summary of cases}
%=======================================
We can give a short summary of the result via a table. Here $q$ stands for a one-dimensional solution and $f$ to its associated function such that $B(fq) = 0$.
 \begin{table}[H]
    \caption{Easy 2D cases}
    \label{tbl:table2d}
    \centering
    \begin{tabular}{cccc}
    \toprule
    Case  & $Q$ & $F(a_1,a_2,t)$\\
    \midrule
    1.1  & 1 & 1 \\
    1.2, 1.3  & $q(x_2,a_2,a_1t)$ & $f(a_2,a_1t)$ \\
    2.1, 3.1 &  $q(x_1,a_1,a_2t)$ & $f(a_1,a_2t)$\\
    2.2  & $\dfrac{\ol a_1 \ol a_2+ \tfrac{c_1 f_1 + c_2 f_2- c_1 c_2}{1-a_1a_2 \delta t}}{1-tS}$ & \\
    2.3  & &\\
    \bottomrule
    \end{tabular}
\end{table}

\endinput
